\chapter{Revisão da Literatura}
\label{cha:revisao}

\section{Metaverso}
\label{sec:metaverso}

O termo metaverso não possui uma única definição consolidada. Ritterbusch e Teichmann (2023), em uma revisão sistemática dedicada especificamente à definição do conceito, identificaram diferentes formulações na literatura e destacaram atributos como ambientes tridimensionais, representação dos usuários por avatares, interação, imersão, economia virtual e tecnologias como realidade virtual, realidade aumentada e blockchain.

A definição adotada neste trabalho acrescenta à dimensão tridimensional e imersiva três características particularmente relevantes para o escopo do projeto: persistência, interoperabilidade e economia digital. A persistência corresponde à continuidade do ambiente mesmo quando determinado usuário deixa a plataforma. A interoperabilidade representa a possibilidade, ainda conceitual ou parcial, de transportar identidades, objetos ou outros elementos entre diferentes ambientes. A economia digital corresponde à existência de mecanismos de troca, propriedade, produção ou consumo de bens e serviços digitais.

A literatura mais recente reforça que o metaverso deve ser compreendido como um ecossistema tecnológico, e não somente como um ambiente visual tridimensional. Tukur et al. (2024) realizaram uma revisão de escopo de trabalhos publicados entre 2020 e 2022 e identificaram quatro técnicas principais e doze classes tecnológicas relacionadas à criação de ambientes digitais para o metaverso. Entre as tecnologias mais recorrentes estavam a realidade estendida, a inteligência artificial e tecnologias descentralizadas.

O estudo de Tukur et al. (2024) também evidencia que as aplicações do metaverso ultrapassam o entretenimento. Os autores identificaram aplicações em onze setores, com destaque para educação, manufatura, saúde e mercado imobiliário. Esse resultado é relevante porque desloca a discussão de uma visão exclusivamente associada a jogos para uma perspectiva de infraestrutura digital e interação entre pessoas, objetos e serviços.

Uddin et al. (2024) ampliam essa perspectiva ao analisar a convergência entre metaverso, inteligência artificial e blockchain. O estudo discute como essas tecnologias podem atuar conjuntamente na criação de ambientes virtuais, na personalização de experiências, na gestão de ativos digitais e em questões de segurança. Essa abordagem ajuda a separar o conceito de metaverso das tecnologias que podem viabilizá-lo: realidade virtual, inteligência artificial e blockchain são componentes possíveis, mas não devem ser tratados automaticamente como sinônimos do conceito.

Assim, a revisão indica uma distinção importante. Algumas propriedades aparecem como características conceituais recorrentes, como ambientes digitais tridimensionais, interação por avatares e experiências compartilhadas. Outras propriedades, como blockchain, inteligência artificial, realidade virtual e economias baseadas em ativos digitais, funcionam principalmente como tecnologias ou mecanismos de implementação. A definição adotada neste trabalho deve preservar essa distinção.

\section{Questões da Pesquisa}
\label{cha:questoes}

Para orientar a análise comparativa, foram formuladas as seguintes questões de pesquisa:

\begin{itemize}
    \item \textbf{QP1:} Quais características aparecem com maior recorrência nas definições científicas de metaverso?
    \item \textbf{QP2:} Quais características da definição adotada neste projeto são sustentadas pela literatura analisada?
    \item \textbf{QP3:} Quais características podem ser unificadas sem ampliar excessivamente o conceito de metaverso?
    \item \textbf{QP4:} Quais características devem ser tratadas como possibilidades de implementação, e não como requisitos necessários?
    \item \textbf{QP5:} Quais lacunas permanecem na definição adotada e quais consequências essas lacunas podem trazer para o desenvolvimento do projeto?
\end{itemize}

As questões foram formuladas para evitar que a comparação se limite a uma escolha entre duas definições. O objetivo é identificar convergências, divergências e níveis diferentes de abstração. Dessa maneira, uma característica pode ser mantida na definição final como elemento conceitual, enquanto outra pode ser deslocada para a camada tecnológica ou para a descrição das aplicações.

\section{Trabalhos Relacionados}
\label{cha:relacionados}

O primeiro trabalho utilizado para comparação é o estudo de Ritterbusch e Teichmann (2023), que realiza uma revisão sistemática especificamente voltada à definição do metaverso. Sua principal contribuição para este projeto é conceitual: o estudo demonstra que não existe consenso absoluto e organiza atributos recorrentes das definições encontradas.

O segundo trabalho, utilizado como referência principal para a comparação aprofundada, é a revisão de escopo de Tukur et al. (2024). Diferentemente do foco predominantemente conceitual de Ritterbusch e Teichmann, Tukur et al. concentram-se nas técnicas, tecnologias e aplicações necessárias à criação de ambientes digitais associados ao metaverso. O estudo utilizou procedimentos de revisão baseados no PRISMA-ScR e pesquisou ACM, IEEE, Scopus e Google Scholar, mantendo trinta trabalhos primários considerados relevantes.

Para este projeto, a comparação central é entre a definição de trabalho adotada e a perspectiva apresentada por Tukur et al. (2024). A definição do projeto enfatiza a experiência do usuário e as propriedades esperadas do ambiente; Tukur et al. enfatizam a infraestrutura tecnológica e os setores de aplicação. Portanto, os dois enfoques não são concorrentes: eles descrevem dimensões diferentes do mesmo objeto.

\begin{table}[H]
\centering
\caption{Comparação entre a definição do projeto e o estudo de Tukur et al. (2024)}
\label{tab:comparacao}
\begin{tabular}{p{0.24\textwidth}p{0.33\textwidth}p{0.33\textwidth}}
\toprule
\textbf{Dimensão} & \textbf{Definição do projeto} & \textbf{Tukur et al. (2024)}\\
\midrule
Ambiente tridimensional & Elemento central da definição & Relacionado à criação de ambientes virtuais\\
Imersão & Destacada na experiência do usuário & Associada principalmente à XR e à convergência de VR\\
Avatares & Meio de representação e interação & Compatível com a interação em ambientes virtuais\\
Persistência & Considerada característica desejável do ambiente & Não aparece como eixo principal da taxonomia tecnológica\\
Interoperabilidade & Considerada como possibilidade estrutural & Relacionada aos desafios de integração e conexão entre tecnologias\\
Economia digital & Considerada parte da experiência ampliada & Aparece como dimensão de aplicação e como consequência do ecossistema tecnológico\\
IA & Tecnologia de suporte & Uma das tecnologias mais recorrentes\\
Blockchain/descentralização & Tecnologia de suporte à economia e aos ativos & Uma das tecnologias recorrentes\\
Aplicações & Reuniões, aulas, shows, compras e experiências digitais & Onze setores identificados na revisão\\
\bottomrule
\end{tabular}
\end{table}

A comparação mostra que não há necessidade de escolher entre as duas perspectivas. A definição do projeto pode funcionar como uma camada conceitual orientada à experiência e às propriedades do ambiente, enquanto a revisão de Tukur et al. fornece a camada tecnológica necessária para explicar como tais propriedades podem ser implementadas. Essa complementaridade fundamenta a proposta de unificação apresentada posteriormente.
